\section{问题重述}
题目计划在东经$98.5^\circ$、北纬$39.4^\circ$、海拔$3000$ m处建设半径为$350$ m的圆形定日镜场。镜场坐标系以建设区域中心为原点，$x$轴向东、$y$轴向北、$z$轴竖直向上。规划吸收塔高度为$80$ m，塔顶集热器为高$8$ m、直径$7$ m的圆柱形外表受光式集热器，吸收塔周围$100$ m内不得布置定日镜。

定日镜为宽度不小于高度的矩形，边长均须位于$2$--$8$ m，安装高度位于$2$--$6$ m，并保证镜面转动时不触地。相邻定日镜底座中心距离应比镜面宽度至少多$5$ m。题目规定年均指标使用每月21日$9{:}00$、$10{:}30$、$12{:}00$、$13{:}30$和$15{:}00$五个时点计算。

\textbf{问题一：}吸收塔位于区域中心，全部定日镜尺寸为$6$ m$\times6$ m、安装高度为$4$ m，镜心位置由附件给定。计算月平均和年平均光学效率、输出热功率及单位镜面面积输出热功率。

\textbf{问题二：}在全部定日镜尺寸及安装高度相同的条件下，设计吸收塔位置、镜面尺寸、安装高度、镜数和镜位，使年平均输出热功率达到$60$ MW，并使单位镜面面积年平均输出热功率尽量大。

\textbf{问题三：}允许各定日镜尺寸和安装高度不同，在相同额定功率条件下重新设计镜场，并继续提高单位镜面面积年平均输出热功率。
